\documentclass[aspectratio=169,notheorems]{corvinusmetropolis}

\usepackage[hungarian]{babel}
\usepackage[style=apa, backend=biber]{biblatex}
\addbibresource{05-05.bib}
\usepackage{amsmath}

\renewcommand{\myinstlogo}{./corvinus-logos/black_simple_HUN.pdf}

\title{Biodiversity \& Nature  in Europe}
\subtitle{Biodiverzitás és Nature kockázat megjelenése az európai piacon, elmélet és piaci realitás}
\author{Molnár Marcell}
\email{marcell.molnar@stud.uni-corvinus.hu}

\institute[Corvinus University]
  {
  Pénzügy és számvitel alapszak\\
  Pénzügy specializáció\\
  Email: \putemail\\
  }
\date{\today}
\titlegraphic{\hfill{\includegraphics[width=.2\textwidth]{\myinstlogo}}}
\keywords{nature_finance, biodiversity_finance, sustainable_finance}

\begin{document}
\openingframes{}

\begin{frame}
  \frametitle{Irodalomi háttér}
  Az elmúlt években a Fenntartható Pénzügyek kutatói fokozott figyelmet fordított az új publikációkban a Biodiversity Finance ágnak\\ 
  Friss, releváns kutatások:
  \begin{itemize}
    \item \alert{biodiverzitási prémium} jelenléte és mértéke \parencite{coqueret_biodiversity_2024, naffa_biodiversity_2024},
    \item \alert{ENCORE adatbázis} biodiverzitás témakörben \parencite{clement_mapping_2025},
    \item \alert{private capital befektetések} \parencite{flammer_biodiversity_2025}\\
  \end{itemize}
  \textbf{A hozzájárulásunk a szakirodalomhoz a biodiversity/nature hatások empirikus vizsgálata az európai részvénypiacon}
\end{frame}

\begin{frame}
  \frametitle{Kutatási kérdések}
  \begin{enumerate}
    \item A biodiverzitási adatok beépítésével növelni tudjuk a maximálisan elérhető Sharpe-rátát, az egyszerű kiszűrési stratégiához képest? \parencite{pedersen_responsible_2021} 
    \item A piaci hozamokban megjelenik-e a biodiverzitási prémium \parencite{coqueret_biodiversity_2024, naffa_biodiversity_2024}, és ha igen, akkor ennek a méréke mekkora az európai piacon?
  \end{enumerate}
\end{frame}

\begin{frame}
  \frametitle{Módszertan}
  Egy olyan befektetőt feltételezünk (shortolás megengedett), akinek a hasznossági függvényébe a vagyon maximalizálása mellett, \alert{pozitívan gondolkodik a magasabb biodiversity/nature értékekről}.\\
  \textcite{pedersen_responsible_2021} alapján a \alert{hasznossági fv-e lineáris és a következő alakú}:
  $$
  U = E[\widehat{W}|s] - \frac{\hat{\gamma}}{2}Var(\widehat{W}|s) + Wf(\bar{s}),
  $$
  ahol $\widehat{W}=W(1+r^f+x'r)$, $\hat{\gamma}$ az abszolút kockázatelutasító paraméter, $s, \bar{s}$ pedig a biodiverzitási/nature score és számtani átlaga.
\end{frame}

\begin{frame}
  \frametitle{Módszertan}
  Egy sztenderd mean-variance hasznosságfüggvényű befektető a befektetési univerzumban megkeresi az \alert{érintési portfóliót}, majd ezt keveri a \alert{kockázatmentes eszközzel} az egyéni kockázatkerülésének függvényében.
  \\Az érintési portfólió az a portfólió, ami a legmagasabb Sharpe-rátával, ezért \textcite{pedersen_responsible_2021} \textbf{általánosította} ezt azzal, hogy \alert{minden egyes ESG-proxy mellett} (jelenesteben biodiversity score), \alert{meghatározható a maximálisan elérhető Sharpe-ráta}.\\
  \bigskip
  \textbf{A probléma idáig nem függ az egyéni hasznosságtól, azaz a befektetési univerzumra egyértelműen meghatározható ez az ún. BIODIV-SR határ!}
\end{frame}

\begin{frame}
  \frametitle{Adatbázisok, befektetési univerzum}
\end{frame}

\begin{frame}
  \frametitle{Roadmap}
\end{frame}

\printbibliography

% \closingframes{Köszönöm a figyelmet!}
\end{document}
